\documentclass[border=10pt]{standalone}
\usepackage{fontspec}
\setmainfont{Noto Serif}[Ligatures=TeXOff]
\XeTeXgenerateactualtext=1
\newfontfamily\bengalifont{Noto Serif Bengali}[Script=Bengali, NFSSFamily=bnfam]
\usepackage[Latin,Bengali]{ucharclasses}
% group-free font switching (as in header-common.tex): robust inside braces, cells and quotes
\makeatletter
\def\BN@saved{\rmdefault}
\def\BN@bn{bnfam}
\newcommand\BNin{\edef\BN@tmp{\f@family}\ifx\BN@tmp\BN@bn\else\let\BN@saved\BN@tmp\fi\fontfamily{bnfam}\selectfont}
\newcommand\BNout{\fontfamily{\BN@saved}\selectfont}
\makeatother
\setTransitionsFor{Bengali}{\BNin}{\BNout}
% Typographic quotes and dashes are not in a ucharclasses class (class 0), so a Bangla word right
% after “ or ‘ would not trigger the Bangla font. Treat them as Latin so that every switch fires.
\XeTeXcharclass"2018=\XeTeXcharclass`A
\XeTeXcharclass"2019=\XeTeXcharclass`A
\XeTeXcharclass"201C=\XeTeXcharclass`A
\XeTeXcharclass"201D=\XeTeXcharclass`A
\XeTeXcharclass"2013=\XeTeXcharclass`A
\XeTeXcharclass"2014=\XeTeXcharclass`A
\XeTeXcharclass"2026=\XeTeXcharclass`A
\usepackage{newunicodechar}
\newunicodechar{।}{{\bengalifont\char"0964}}
\newunicodechar{॥}{{\bengalifont\char"0965}}

\usepackage[table]{xcolor}
\usepackage{booktabs,tabularx,array}
\definecolor{bapilrule}{HTML}{1F2937}
\definecolor{bapilstripe}{HTML}{F3F4F6}
\renewcommand{\arraystretch}{1.28}
\setlength{\tabcolsep}{5pt}

\begin{document}
\begin{minipage}{17cm}
{\large\bfseries সারণি ৬০।}\ {\large প্রধান দাবিগুলোর প্রমাণগত অবস্থা (২১ ও ২২ অক্টোবর ২০২৬-এ যাচাইকৃত)}\par\vspace{7pt}
\small
\rowcolors{2}{bapilstripe}{white}
\begin{tabularx}{\linewidth}{@{}>{\hsize=1.275\hsize\raggedright\arraybackslash}X>{\hsize=0.637\hsize\raggedright\arraybackslash}X>{\hsize=0.735\hsize\raggedright\arraybackslash}X>{\hsize=1.275\hsize\raggedright\arraybackslash}X>{\hsize=1.078\hsize\raggedright\arraybackslash}X@{}}
\arrayrulecolor{bapilrule}\toprule
\rowcolor{white}\textbf{লেখায় যেভাবে দাবি করা হয়েছে} & \textbf{কোথায়} & \textbf{অবস্থা} & \textbf{সর্বোত্তম প্রাপ্ত প্রমাণ} & \textbf{লেখায় উপস্থাপন} \\
\midrule
জাতিসংঘের EGDI ২০২৪-এ বাংলাদেশ ১৯৩ দেশের মধ্যে ১০০তম, মান ০.৬৫৭ (২০২২ সালে ১১১তম) & অধ্যায় ২; সারণি ৪ & যাচাইকৃত তথ্য & জাতিসংঘের ই-সরকার জ্ঞানভান্ডারের দেশভিত্তিক তথ্য [26]; [5, 21] & প্রাথমিক উৎসসহ উল্লিখিত \\
আইটিইউ গ্লোবাল সাইবার সিকিউরিটি ইনডেক্স ২০২৪-এ বাংলাদেশ টিয়ার ১ (রোল-মডেল) দেশ হিসেবে চিহ্নিত & অধ্যায় ২; সারণি ৪ & যাচাইকৃত তথ্য (সরকারি জাতীয় উৎস) & বিজিডি ই-গভ সার্ট [29] & জাতীয় সার্টের প্রতিবেদন অনুযায়ী উল্লিখিত \\
২০২৩ সালের সরকারি এআই প্রস্তুতি সূচকে ৮২তম & অধ্যায় ২; সারণি ৪ & যাচাইকৃত তথ্য & সূচকের মূল প্রতিবেদন [24]: র‍্যাংক ৮২, মোট স্কোর ৪৬.০৪; সুলতানা ও সহলেখকও [25] উল্লেখ করেছেন & উৎসসহ তথ্য হিসেবে উল্লিখিত \\
প্রতিবেদন অনুযায়ী এপ্রিল ২০২৬-এ ব্যক্তিগত উপাত্ত সুরক্ষা আইন, ২০২৬ প্রণীত & অধ্যায় ১, ২, ১৪; সারণি ৪, ৪২ & উৎস-প্রতিবেদিত তথ্য, স্বাধীনভাবে অযাচাইকৃত & সিকিউরিটি, জুন ২০২৬ [7]; আজম, সেপ্টেম্বর ২০২৬ [38]; সরকারি আইন-ডেটাবেসে শিরোনাম তালিকাভুক্ত (পাঠ পর্যালোচিত নয়) & প্রণয়ন প্রতিবেদিত; কার্যকারিতা নিশ্চিতকরণসাপেক্ষ হিসেবে চিহ্নিত \\
প্রতিবেদন অনুযায়ী ১০ এপ্রিল ২০২৬-এ জাতীয় উপাত্ত ব্যবস্থাপনা বিল, ২০২৬ পাস & অধ্যায় ২, ১৪; সারণি ৪২ & উৎস-প্রতিবেদিত তথ্য & ডেইলি সান, ১০ এপ্রিল ২০২৬ [8] & প্রতিবেদিত; বর্ণিত বৈশিষ্ট্যগুলো ২০২৫ সালের অধ্যাদেশের \\
২০২৫ সালের অধ্যাদেশে প্রধান ডেটা কর্মকর্তা, অভিযোগ ও অপরাধ-সংক্রান্ত ধারা তাৎক্ষণিকভাবে কার্যকর হয়নি & অধ্যায় ৮, ১৪; সারণি ৪২ & উৎস-প্রতিবেদিত তথ্য (মাধ্যমিক আইনি ভাষ্য) & গ্রিনলিফ ও আল রাজী, জুন ২০২৬ [39]; ১৮ মাসের সম্মতি-সময়সীমা প্রতিবেদিত [97] & অধ্যাদেশের ক্ষেত্রে উল্লিখিত; আইনের অধীনে অবস্থান আইনজীবীর পর্যালোচনার জন্য রাখা \\
প্রতিবেদন অনুযায়ী ১০ এপ্রিল ২০২৬-এ সাইবার নিরাপত্তা বিল, ২০২৬ পাস; ৩০ জুন ২০২৬-এ সাইবার নিরাপত্তা (সংশোধন) বিল, ২০২৬ পাস, যার বিষয়বস্তু ভিন্নভাবে প্রতিবেদিত; জুলাই ২০২৬ থেকে আরেকটি খসড়া সংশোধনী পর্যালোচনাধীন & অধ্যায় ২, ১৪; সারণি ৪২ & উৎস-প্রতিবেদিত তথ্য; সংশোধনীর বিষয়বস্তু নিয়ে উৎসগুলো ভিন্নমত & [8], [9], [40], [41] & প্রতিবেদিত; বিষয়বস্তু ও কার্যকারিতা নিশ্চিতকরণসাপেক্ষ \\
জানুয়ারি ২০২৬-এ খসড়া জাতীয় এআই নীতি ২০২৬-২০৩০ জন-পরামর্শে প্রকাশিত; ৯ ফেব্রুয়ারি ২০২৬ তারিখের সংশোধিত খসড়া সংস্করণ ২.০ & অধ্যায় ১, ২, ১৪ & উৎস-প্রতিবেদিত তথ্য (জানুয়ারির খসড়া); সরকারি পোর্টালের তালিকা (সংস্করণ ২.০; বিষয়বস্তু পর্যালোচিত হয়নি) & দ্য ডেইলি স্টার, ২৬ জানুয়ারি ২০২৬ [10]; দ্য ফাইন্যান্সিয়াল এক্সপ্রেস, ২৯ জানুয়ারি ২০২৬ [11]; aipolicy.gov.bd [12] & প্রতিবেদিত \\
খসড়া নীতি এআই-সৃষ্ট মিথ্যা বা বিভ্রান্তিকর কনটেন্ট মোকাবিলার কথা বলে; এআই-সৃষ্ট লেখায় লেবেল BAPIL-এর নিজস্ব নিয়ম & অনুচ্ছেদ ১৩.৮; সারণি ৪১; সংযুক্তি ঢ & উৎস-প্রতিবেদিত তথ্য (নীতির বিষয়বস্তু); নকশাগত সুপারিশ (লেবেল) & [10] ও [11]-এ এআই-সৃষ্ট মিথ্যা বা বিভ্রান্তিকর কনটেন্ট প্রতিরোধের ব্যবস্থার বর্ণনা আছে; কোনোটিতেই লেবেলের বাধ্যবাধকতার উল্লেখ নেই & লেবেলের বিধান খসড়া নীতির বলে উল্লেখ করা হয়নি \\
একটি যৌথ দলের মাধ্যমে ২০২৬-এর চতুর্থ ত্রৈমাসিকে জাতীয় এআই নীতি প্রণয়ন শুরু & অধ্যায় ১, ১৪, ২১ & উৎস-প্রতিবেদিত তথ্য & দ্য ডেইলি অবজারভার, ১০ আগস্ট ২০২৬ [14] & প্রতিবেদিত \\
১৯২.৬৬ কোটি টাকার এআই হাব কারিগরি সহায়তা প্রকল্প জুন ২০২৬-এ অনুমোদনের অপেক্ষায়; আগস্ট ২০২৬-এ ১৩ মিলিয়ন মার্কিন ডলারের অনুদান-চুক্তি স্বাক্ষরিত & অধ্যায় ৩, ১৭; সারণি ৪২ & উৎস-প্রতিবেদিত তথ্য & [53]; [55] & শুধু মাত্রা বোঝাতে ব্যবহৃত; অনুমোদিত বরাদ্দ হিসেবে বর্ণিত নয় \\
রেফারেন্স বাস্তবায়ন গিটহাব রিলিজ v1.1.0-এ প্রকাশিত এবং ব্লুপ্রিন্টের Zenodo রেকর্ডের সঙ্গে সংরক্ষিত & সূচনা-অংশ; অধ্যায় ১, ১১, ১৬, ২১; সারণি ১ & রচনার সময় ভবিষ্যৎ মাইলফলক & রিপোজিটরিতে ৪৯টি উত্তীর্ণ ইউনিট পরীক্ষাসহ কোড বিদ্যমান; ২২ অক্টোবর ২০২৬ পর্যন্ত রিলিজ ও ব্লুপ্রিন্টের DOI প্রকাশ্য হয়নি & রিলিজটি কীভাবে যাচাই করা যায় তা পাঠককে জানানো হয়েছে; রিলিজ প্রকাশের পরই DOI রেকর্ড প্রকাশিত হবে (সংযুক্তি ঘ) \\
প্রদর্শনীতে Recall@10 ০.৮৮ (২৪টি উত্তরযোগ্য প্রশ্নের মধ্যে ২১টি) এবং ৬টির মধ্যে ৬টি সঠিক বিরতি & অনুচ্ছেদ ১১.৫ & গবেষণা-ফলাফল (ছোট প্রশ্নসেটে উন্নয়নকালীন ফলাফল) & বিল্ডে তৈরি মূল্যায়ন প্রতিবেদন: ৮টি নথি, প্রায় ১,৭০০টি অনুচ্ছেদ, ২৪টি প্রশ্ন & সীমাবদ্ধতাসহ উন্নয়নকালীন ফলাফল হিসেবে উপস্থাপিত \\
Recall@10 কমপক্ষে ০.৮৫ ও উদ্ধৃতির নির্ভুলতা কমপক্ষে ০.৯৫ & অধ্যায় ৯, ১২; সারণি ২৮, ৩৮ & প্রস্তাবিত লক্ষ্য & এখনো পরিমাপ করা হয়নি & লক্ষ্য হিসেবে চিহ্নিত; পরিমাপ-প্রটোকল সংযুক্তি দ-তে \\
৩৬ মাসে প্রায় ২০.৪ কোটি টাকার নির্দেশক বাজেট (মূল হিসাব) & অধ্যায় ১৭; সারণি ৪৯, ৫০, ৬১, ৬২; সারসংক্ষেপ & পরিকল্পনা-অনুমান (লেখকের বিশ্লেষণ) & সারণি ৪৯-এর অনুমান থেকে blueprint/\allowbreak{}src/\allowbreak{}budget\_model.py দ্বারা হিসাবকৃত; আসন্নীকরণের আগের মোট ২০.৪৩ কোটি টাকা & নির্দেশক; পরিস্থিতি, সংবেদনশীলতা ও জানা বাদ-পড়া ব্যয় উল্লিখিত (অনুচ্ছেদ ১৭.২; সংযুক্তি থ) \\
বছর ৪ থেকে ৮-এ প্রায় ২১৮.৯ কোটি টাকার শর্তসাপেক্ষ জাতীয় সম্প্রসারণ, বছর ৮-এর মধ্যে প্রায় ১১৮টি প্রত্যক্ষ পদ & অনুচ্ছেদ ১৭.৩; সংযুক্তি ব; সারণি ৬৯, ৭০; সারসংক্ষেপ & পরিকল্পনা-অনুমান (লেখকের বিশ্লেষণ) & সারণি ৪৯-এর একক-ব্যয় দিয়ে blueprint/\allowbreak{}src/\allowbreak{}budget\_model.py (scaleup ও jobs ফাংশন) দ্বারা হিসাবকৃত & শর্তসাপেক্ষ ও অনুরোধবহির্ভূত হিসেবে উপস্থাপিত; ব্লুপ্রিন্ট ২.০-এ একটি ব্যয়-নির্ধারিত পরিকল্পনা এর স্থলাভিষিক্ত হবে (U17) \\
নবম জাতীয় বেতন স্কেলে গ্রেড ৯-এর মূল বেতন ৪৪,০০০ টাকা & অনুচ্ছেদ ১৭.৫; সারণি ৪৯ & উৎস-প্রতিবেদিত তথ্য (মন্ত্রিসভার অনুমোদন) & দ্য ফাইন্যান্সিয়াল এক্সপ্রেস, ৩১ আগস্ট ২০২৬ [100]; গেজেট দেখা হয়নি & অবস্থাসহ উল্লিখিত \\
১৯ অক্টোবর ২০২৬-এ বিনিময় হার প্রতি মার্কিন ডলারে ১২৩.৬৫ টাকা & সারণি ৪৯ & যাচাইকৃত তথ্য & বাংলাদেশ ব্যাংক [107] & ১২৪ টাকার পরিকল্পনা-হারের ভিত্তি \\
সেপ্টেম্বর ২০২৬-এ মূল্যস্ফীতি ৮.৩৪\% & সারণি ৪৯ & উৎস-প্রতিবেদিত তথ্য & ঢাকা ট্রিবিউন, ৭ অক্টোবর ২০২৬, বিবিএস-এর তথ্য উদ্ধৃত করে [108] & ৮\% মূল্যবৃদ্ধি-অনুমানের ভিত্তি \\
একটি L40S GPU কার্ড প্রায় ৭,৫০০ মার্কিন ডলার & সারণি ৪৯ & উৎস-প্রতিবেদিত তথ্য (বিক্রেতার প্রকাশনা) & মোডাল, আগস্ট ২০২৫ [109] & শুধু মূল্যের ভিত্তি; দরপত্র প্রয়োজন \\
বিশ্লেষকদের সাশ্রয়কৃত সময়ের মূল্য বছরে প্রায় ০.৭ কোটি টাকা & অনুচ্ছেদ ১৭.৫ & লেখকের বিশ্লেষণ & [100] ও অনুমিত ব্যবহারের ভিত্তিতে হিসাব & দৃষ্টান্তমূলক হিসেবে চিহ্নিত \\
উল্লিখিত ক্ষমতাসহ স্টিয়ারিং বোর্ড, ERRC ও TAB & অধ্যায় ৫; সংযুক্তি ঙ থেকে ছ & নকশাগত সুপারিশ & কোনো সংস্থা গঠিত বা অনুমোদিত হয়নি & প্রস্তাব হিসেবে উপস্থাপিত; অবস্থা, নিয়োগ ও সীমা সংযুক্তি ধ-তে \\
২০২৮ থেকে ২০৪১ পর্যন্ত দীর্ঘমেয়াদি লক্ষ্য & অধ্যায় ১৬; সারণি ৪৬ & প্রস্তাবিত লক্ষ্য & প্রযোজ্য নয় & নির্দেশক; কৌশলগত পর্যালোচনায় সংশোধিত হবে \\
\bottomrule
\end{tabularx}
\par\vspace{5pt}{\footnotesize\itshape অবস্থা-নির্দেশক লেবেলগুলোর সংজ্ঞা সংযুক্তি ত-তে দেওয়া আছে। ওয়েব উৎসগুলো ২১ ও ২২ অক্টোবর ২০২৬-এ দেখা হয়েছে। সরকারি আইন-ডেটাবেস স্বয়ংক্রিয় সংগ্রহের অনুমতি দেয় না, তাই আইনের পাঠ সরাসরি পর্যালোচনা করা হয়নি।}
\par\vspace{3pt}{\scriptsize\color{gray} উৎস: BAPIL টেকনিক্যাল ব্লুপ্রিন্ট, সংস্করণ ১.০ (২০২৬)। CC BY 4.0।}
\end{minipage}
\end{document}
